

In a parametric model, the maximum of divergence happens for a prior distribution
that is concentrated at a single value of the parameter $\paramValue$.
That means, for a parametric model, which allows any high peak prior, instead of maximizing over the set of
all acceptable distributions, we can simply maximize over the parameter space, using the conditional version
of the divergence measure.

\begin{definition}[Generic Parametric Model]
    We say a parametric model \pmodel is \emph{generic}, if for any arbitrary probability distribution $r$,
    there is a probability measure $R \in \set{P}$, such that for every $\rval[]{s} \in \set{s}$
    \[
        \marginalEq[r]{\rval[]{s}}\cdot
    \]
\end{definition}

\begin{lemma}
    \label{lm:DxP_uppper_bound}
    Assume that \pmodel is a parametric model.
    Let \Xpairs be a query set in \Mt.
    For every $P \in \set{P}$ we have
    \begin{equation*}
        \Dx \leq \supParam \DxCond
    \end{equation*}
\end{lemma}

\begin{proof}
    \newcommand{\RP}{\sumixe \pi(\xeIt[i]) \ln \frac{R(\xIe[i])}{P(\xIe[i])}}

    Let $R$ and $P$ be two arbitrary probability measures in \Mt.
    For any parameter value \paramValue we have:
    \[
        \DxCond[P] - \DxCond[R] = \RP \cdot
    \]
    The conditional divergence is nonnegative and we have
    \[
        \RP \leq \DxCond \cdot
    \]
    This inequality holds for every \paramValue.
    Therefore, the supremum of the right-hand side w.r.t. \paramValue is also greater than or equal to the
    supremum of the left-hand side:
    \begin{equation}
        \supParam \RP \leq \supParam \DxCond \cdot
        \label{eq:ineq-theorem1}
    \end{equation}
    Trivially, for any \paramValue we have
    \[
        \RP \leq \supParam \RP \cdot
    \]
    \Mt is a parametric model and $\marginalEq[R]{\xe[i]}$, where $r(\paramValue)$ is
    a prior probability distribution.
    Since the above inequality holds for every value of \paramValue and $r(\paramValue) \geq 0$, we can multiply both
    sides by $r(\paramValue)$, and integrate.
    Notice that the right hand side is not a function of $\paramValue$ and it remains unchanged
    because $\int \dt[r] = 1$:
    \[
        \DxExpand[R] \leq \supParam \RP\cdot
    \]
    Now from Inequality~\ref{eq:ineq-theorem1}, for \emph{any} $P \in \set{P}$ and $R \in \set{P}$, it
    follows that
    \[
        \DxExpand[R] \leq \supParam \DxCond,
    \]
    which proves the lemma.
\end{proof}

\begin{theorem}
    \label{th:Dx_parametric_form}
    Suppose \pmodel is a parametric model, and let $\Xpairs$ be a query set in \Mt.
    If for every $\epsilon > 0$ and every value of the parameter $\paramValue$, there exists a distribution $R \in \set{P}$
    such that
    \begin{equation}
        \DxCond[][\paramValue] - \DxExpand[R]  < \epsilon,
        \label{eq:epsilon_close}
    \end{equation}
    then for any $P \in \set{P}$ we have:
    \begin{equation}
        \Dx = \supParam \DxCond  \cdot
        \label{eq:div_parametric}
    \end{equation}
    \label{th:div_for_parametric}
\end{theorem}

\begin{proof}
    Assume, for the sake of contradiction, that
    \[
        \Dx \neq \supParam \DxCond \cdot
    \]
    Then, Lemma~\ref{lm:DxP_uppper_bound} implies $\Dx < \supParam \DxCond$.
    Supremum is the minimal upper bound of a set.
    That means $\Dx$ can not be an upper bound for the set of values of $\DxCond$,
    and there should exist a $\paramValue^{*}$ such that
    \[
        \DxCond[][\paramValue^{*}] > \Dx \cdot
    \]

    By the theorem assumption, for $\paramValue^{*}$ and $\epsilon = \DxCond[][\paramValue^{*}] - \Dx$ there should be a
    distribution $R \in \set{P}$ such that Equation~\ref{eq:epsilon_close} holds, and
    \begin{align*}
        \DxExpand[R] > \Dx  \cdot
    \end{align*}
    On the other hand
    \[
        \supQ \DxExpand[Q] \geq \DxExpand[R] \cdot
    \]
    Notice that the left hand side is $\Dx$.
    This contradiction implies that the assumption can not hold.
    Hence,
    \[
        \Dx = \supParam \DxCond. \qedhere
    \]
\end{proof}

\begin{proposition}
    A generic parametric model satisfies the conditions of Theorem~\ref{th:Dx_parametric_form}.
\end{proposition}


A vector in $\mathbb {R}^{n}$ shares many properties with a real function.
We sometimes assume that a real function is a real vector, and vice versa.
In the next theorem we will use the minimax theorem, assuming $x$ is a real function.



\newcommand{\Pentry}{p_{\xIe[i]}}
\newcommand{\PPentry}{\sad{\Pentry}}
\newcommand{\sad}[1]{#1^{*}}
\newcommand{\rDist}{r}
\newcommand{\fExpand}{\intProd[\r]{\left( \fInside{\Pentry} \right)}}
\newcommand{\fInside}[1]{\sumixe \pi(\xIet[i]) \ln \frac{\pi(\xIet[i])}{#1}}

\newcommand{\RR}{\sad{R}}
\newcommand{\Rvector}{\sad{R}_\Qset}
\newcommand{\RRvector}{\mathbf{\sad{R}_\Qset}}
\newcommand{\Cvector}{\mathbf{\central}}

\newcommandx{\f}[2][1 = \r, 2 = \P, usedefault]{ f \ifthenelse{\isempty{#1}}{}{(#2, #1)}}
\renewcommand{\r}{r}
\renewcommand{\P}{\mathbf{P_{\Qset}}}
\newcommand{\rr}{\sad{\r}}
\newcommand{\PP}{\mathbf{\sad{P}_{\Qset}}}
\newcommand{\R}[1][{\xIe[i]}]{\sad{R}({#1})}
\newcommand{\saddle}{(\PP, \rr)}


\begin{theorem}
    Let \pmodel be a generic parametric model, $\central \in \set{P}$, and \Xpairs be a query set.
    Suppose that there is a prior distribution $\cMass{}$ where $\cMass{\paramValue} > 0$ for all \paramValue,
    and for every $i$, \x[i] and \e[i] we have
    \[
        \central(\xIe[i]) = \int \pi(\xIet[i]) \dt[\cMass{}] \cdot
    \]
    If there is a constant $\lambda$ such that
    \begin{equation*}
        \DxCond[\central] = \lambda,
    \end{equation*}
    then \central is a central distribution for \Qset.
\end{theorem}

\begin{proof}
    We define
    \begin{equation*}
        \f = \fExpand,
    \end{equation*}
    where $\r$ is a real function of the model parameter \paramValue, and $\P$ is a vector of real
    numbers with an entry $\Pentry$ for every $i, \xe[i]$.

    Let $\saddle$ be a simultaneous solution to the following two optimization problems:

    \begin{equation*}
        \begin{aligned}
            &\text{minimize} & &\f[\rr] & &\text{w.r.t.} \ \P \\
            &\text{subject to} & &\sum_{\x[i]} \Pentry = 1,  & &\text{for all}\  i,\e[i] \\
            \\
            &\text{maximize} & &\f[][\PP] & &\text{w.r.t.} \ \r \\
            &\text{subject to} & &\r(\paramValue) \geq 0, & &\text{for all} \ \paramValue  \\
            & & &\intProd[\r]{} = 1 &
        \end{aligned}
    \end{equation*}
    Clearly, if $\saddle$ exists, it will be a saddle point for $\f[ ]$ on a restricted domain.

    Both of these optimization problems are convex and differentiable~\cite{0521833787}.
    All the constraint functions are affine, and therefore are convex and differentiable.
    Note that for the constraint $\r(\paramValue) \geq 0$, we can define an infinite number of affine
    constraint functions of the form $\int \delta(\paramValue - t) \r(t) dt \geq 0$, where $\delta$ is the Dirac delta
    function.

    Additionally, $\f[ ]$ is differentiable with respect to both $\r$ and $\P$.
    For any fixed $\P$, $\f[\cdot]$ is an affine function and is concave.
    For a fixed solution to the second optimization problem denoted by $\rr$, we can write
    \begin{equation}
        \f[\rr][\P] = c - \sumixe \R[{\xe[i]}] \ln \Pentry,
    \end{equation}
    where $c$ and $\R[{\xe[i]}] = \marginal[\rr]{\xe[i]}$ are not functions of $\P$.
    Since $\rr$ is a solution to the second optimization problem, it is feasible and satisfies
    all the constraints.
    Therefore, we have $\sad{R}(\xe[i]) \geq 0$ for any $i$, \e[i], and $\f[\rr][\cdot]$ is the negative of a
    linear combination of concave $\ln \Pentry$ functions with nonnegative coefficients.
    This function is convex.

    Since both optimization problems are convex and differentiable, any point $\saddle$ that simultaneously satisfies
    the KKT conditions~\cite{0521833787} for both problems is a saddle point for $\f[ ]$ on the restricted domain
    defined by the constraints.
    The following KKT conditions are obtained for every \paramValue, $i$, \x[i], and \e[i]:
    \begin{gather}
        \label{eq:kkt_R_is_pxIe}
        \intProd[\rr]{\pi(\xIet[i])} = \lambda(\e[i]) \PPentry, \\
        \label{eq:kkt_pxIe_sums_one}
        \sum_{\x[i]} \PPentry = 1, \\
        \nonumber
        \fInside{\PPentry} - \mu(\paramValue) = \lambda, \\
        \label{eq:kkt_r_sums_one}
        \intProd[\rr]{} = 1, \\
        \label{eq:kkt_r_nonnegative}
        \rr(\paramValue) \geq 0, \\
        \nonumber
        \mu(\paramValue) \geq 0, \\
        \label{eq:kkt_mu_product}
        \mu(\paramValue) \rr(\paramValue) = 0 \cdot
    \end{gather}

    By summing Equation \ref{eq:kkt_R_is_pxIe} over \x[i] and using Equations~\ref{eq:kkt_r_sums_one}
    and~\ref{eq:kkt_pxIe_sums_one}, we obtain $\lambda(\e[i]) = 1$.

    Note that we can impose more restrictive conditions and still derive a set of sufficient conditions for a saddle
    point.
    Hence, we can replace Inequality~\ref{eq:kkt_r_nonnegative} with the more restrictive $\rr(\paramValue) > 0$
    condition.
    By Equation~\ref{eq:kkt_mu_product}, this implies $\mu(\paramValue) = 0$, and the KKT conditions can be
    simplified to
    \begin{gather}
        \label{eq:kkt_Rlambda}
        \fInside{\R} = \lambda, \\
        \nonumber
        \R = \intProd[\rr]{\pi(\xIet[i])}, \\
        \nonumber
        \rr(\paramValue) > 0 \cdot
    \end{gather}

    Let \pmodel be a parametric model.
    If a probability measure $\RR \in \set{p}$ which is obtained by a positive prior $\rr$,
    satisfies Equation~\ref{eq:kkt_Rlambda}, then the KKT conditions imply that $(\RRvector, \rr)$ is a solution to
    both optimization problems, where $\RRvector$ is the corresponding vector of $\RR(\xIe[i])$ probabilities.

    Suppose \central is a central distribution for \Qset.
    Clearly, the point $(\Cvector, \rr)$ satisfies the constraints of the minimization problem (the first problem).
    Therefore, we have
    \begin{equation}
        \label{eq:kkt_saddle_lq_dx}\f[\rr][\RRvector] \leq \f[\rr][\Cvector] \cdot
    \end{equation}
    Notice that by the definition of $\f[ ]$, we have
    \[
        \f[\rr][\Cvector] = \intProd[\rr]{\DxCond[\central]} \cdot
    \]
    Since $\rr$ is a distribution and integrates to one we have
    \[
        \intProd[\rr]{\DxCond[\central]} \leq \supParam \DxCond[\central] = \Dx[\central] \cdot
    \]
    On the other hand, $\RR$ satisfies Equation~\ref{eq:kkt_Rlambda}, and the definition of the divergence
    measure implies
    \[
        \Dx[\RR] = \supParam \DxCond[\RR] = \lambda = \f[\rr][\RRvector] \cdot
    \]
    Thus,
    \[
        \Dx[\RR] \leq \Dx[\central] \cdot
    \]
    $\Dx[\central]$ is minimal in \pmodel, and $\RR \in \set{P}$ implies
    \[
        \Dx[\RR] = \Dx[\central],
    \]
    Consequently, $\RR$ is a central distribution for \Qset.
\end{proof}

\begin{theorem} [minimax]
    Let \pmodel be a generic parametric model and \Xpairs be a query set.
    Let $\f[ ]$ be the function defined by:
    \begin{equation*}
        \f = \fExpand,
    \end{equation*}
    where $\r$ is a probability distribution, and $\P$ is a vector with an
    entry $\Pentry \geq 0$ for every $i, \xe[i]$, such that $\sum_{\x[i]} \Pentry = 1$.

    If $\f[ ]$ has a saddle point, then the central distribution for \Qset can be obtained by:
    \begin{equation}
        \central = \arg\maxP \int \DxCond \dt
        \label{eq:Cx_by_max}
    \end{equation}
\end{theorem}

\begin{proof}
    Hence, we can use Theorem~\ref{th:minimax}
    \begin{equation}
        \minP \max_{r} \f = \max_{r} \minP \f\cdot
        \label{eq:divergence_minmax}
    \end{equation}
    From Equation~\ref{eq:max_for_P} we can see that for an arbitrary fixed $r$ we have
    \[
        \minP \f = \maxP \sumixe R(\xe[i]) \ln P(\xIe[i])\cdot
    \]

    We know that the Kullback–Leibler divergence, and therefore its expected value, is non-negative.
    Consequently:
    \[
        \DxExpand[R] \geq 0\cdot
    \]
    By simple algebra we get
    \[
        \sumixe R(\xe[i]) \ln P(\xIe[i]) \leq \sumixe R(\xe[i]) \ln R(\xIe[i])\cdot
    \]
    Recall that \Mt is a generic parametric model.
    That implies \inP{R}, and we can reach this upper bound by letting $P = R$.
    Hence, for a fixed $r$
    \[
        \minP \f = \int \DxCondExpand[R] \dt[r] \cdot
    \]
    By Equation~\ref{eq:divergence_minmax}, we conclude
    \[
        \central = \arg\max_{\inP{R}} \int \DxCond[R] \dt[r] \cdot \qedhere
    \]
\end{proof}

Let us now use this theorem for our dice example and rewrite Equation~\ref{eq:joint_form}:
\begin{align*}
    \central &= \arg\maxP \int \KlCond[][\seqRval{n}][ ] \dt \\
    &= \arg\maxP \int \sum_{\seqRval{n}} p(\seqRval{n}, \paramValue) \ln
    \frac{p(\paramValue \mid \seqRval{n})}{p(\paramValue)} d\paramValue \cdot
\end{align*}
This means that the central distribution for our dice example, when we consider a query set equivalent to the
sample space, is the probability measure that is obtained by using the Reference prior~\cite{BERNARDO200517} for
the multinomial parameters.

\newcommand{\ratio}[1]{\frac{\avgPost[{#1}][\cMass{}]}{\cMass{#1}}}
\newcommandx{\ratioFull}[3][1=\w, 2=\w, 3=\cMass{}, usedefault]{
    \frac{\avgPost[{#1}][#3]}{#3(#2)}
    \prod_i \frac{\avgPost[{\w[i]}][#3]}{#3(\w[i])}
}
\newcommandx{\pusFrontal}[2][1=\avgPost, 2=\rho]{
    \left(\frac{#1(\w)}{#2(\w)} \prod_i \frac{#1(\w[i])}{#2(\w[i])} \right)
}
\newcommand{\fixedPoint}{\mu \cMass{}}
\newcommand{\fpMap}[1]{T \ifthenelse{\isempty{#1}}{}{\left[ #1 \right]}}

\begin{theorem}
    \label{th:fixed_point}
    Assume \pmodel is a generic parametric model, and $\Qset = \{\X, \seqRv{k}\}$ is a query set in \Mt.
    Let $\W[i]$, for every $i$, be a function of the parameter $\param$ that fully parametrizes \X[i],
    That is, for $\w[i] = \W[i](\paramValue)$ we have:
    \begin{equation}
        \pi(\xIt[i]) = \pi \left( \xIw[i] \right)\cdot
        \label{eq:func_of_theta}
    \end{equation}
    Also, let $\W(\paramValue) = \vSeq{\w}{k} = \w$.

    We define a map $\fpMap{}$ on the space of unnormalized probability measures of \w as follows:
    \begin{equation}
        \fpMap{\rho}(\w) = \pusFrontal[{\avgPost[][\rho]}][\rho]^{\frac{1}{k+1}} \rho(\w) ,
        \label{eq:pus_frontal}
    \end{equation}
    where $\rho$ is an unnormalized probability measure and:
    \begin{align*}
        \avgPost[{\w[i]}][\rho] &= \exp \left\{\sum_{\x[i]} \pi(\xIw[i]) \ln \rho(\wIx[i]) \right\},\\
        \avgPost[\w][\rho] &= \int_{\omegaArea} \avgPost[\paramValue][\rho] d\paramValue,\\
        \avgPost[\paramValue][\rho] &= \exp \left\{\sum_{\x} \pi(\xIt) \ln \rho(\tIx) \right\} \cdot\\
    \end{align*}

    If $\cMass{}$ denotes the probability mass function corresponding to the central distribution for \Qset in \Mt,
    then $\fixedPoint$ is a fixed point for $\fpMap{}$, where
    \[
        \mu = \exp \left\{\frac{1}{k+1}\int \DxCond[\central] \dt[\cPrior[]] \right\} \cdot
    \]
    Moreover, we have
    \begin{equation}
        \cMass{\tIw} = \frac{\avgPost[\paramValue][\cMass{}]}{\avgPost[\w][\cMass{}]} \cdot
        \label{eq:conditional_for_theta}
    \end{equation}
\end{theorem}

\begin{proof}
    Let $Q^*(y) = \convexCond{Q}{y}$, for some \inP{Q_1, Q_2}, $0 \leq \alpha \leq 1$ and
    an arbitrary random variable \rv{Y}.
    \Mt is a parametric model.
    Therefore:
    \begin{align*}
        Q^*(y) &= \alpha \marginal[q_1]{y} + (1 - \alpha) \marginal[q_2]{y}\\
        &= \int \pi(y \mid \paramValue) (\convexCond{q}{\paramValue}) d\paramValue\cdot
    \end{align*}
    $\convexCond{q}{\paramValue}$ is a valid probability distribution and \Mt is a generic parametric model.
    That implies \inP{Q^*}.
    This proves that the set of distributions of \X and each \X[i] is a convex set.

    Thus, the conditions of Theorem~\ref{th:lambda} holds and the central distribution of \Qset satisfies
    \begin{equation*}
        \DxCond[\central ] = \lambda \cdot \\
    \end{equation*}
    Moreover, from Equation~\ref{eq:func_of_theta} and Bayes' theorem it follows
    \begin{align*}
        \DxCond[\central] &= \KlCond[\central][\rval{x}][ ] + \DxCondExpandFull[\central][][ ] \\
        &= \sum_{\rval{x}} \pi(\xIt) \ln \frac{\cMass{\tIx}}{\cMass{\paramValue}} +
        \sumix \pi(\xIw[i]) \ln \frac{\cMass{\wIx[i]}}{\cMass{\w[i]}}\\
        &= \ln \ratio{\paramValue} + \sum_i \ln \ratio{\w[i]}\cdot
    \end{align*}
    That implies
    \begin{equation}
        \label{eq:fixed_point_proof}
        e^{\lambda} = \ratioFull[\paramValue][\paramValue] \cdot
    \end{equation}
    For $\W(\paramValue) = \w$, we have $\cMass{\paramValue, \w}=\cPrior$.
    Hence,
    \begin{equation*}
        \cPrior = \cMass{\tIw} \cMass{\w} \cdot
    \end{equation*}
    We substitute this equation into Equation~\ref{eq:fixed_point_proof}:
    \begin{equation}
        \cMass{\tIw} e^{\lambda} = \ratioFull[\paramValue][\w] \cdot
        \label{eq:fixed_point_proof_1}
    \end{equation}
    Since this equation holds for any value of \paramValue, we can integrate both sides with respect to \paramValue:
    \begin{equation*}
        \int_{\omegaArea} \cMass{\tIw} e^{\lambda} d\paramValue =
        \int_{\omegaArea} \ratioFull[\paramValue][\w] d\paramValue \cdot
    \end{equation*}
    Therefore,
    \begin{equation}
        e^{\lambda} = \ratioFull \cdot
        \label{eq:fixed_point_proof_2}
    \end{equation}
    Now we evaluate $\fpMap{}$ at $\fixedPoint$:
    \begin{equation*}
        \fpMap{\fixedPoint}(\w) =
        \left( \ratioFull[][][\fixedPoint] \right)^{\frac{1}{k+1}} \fixedPoint(\w) =
        e^{\frac{\lambda}{k+1}} \cMass{\w} \cdot
    \end{equation*}
    Clearly for $\mu = e^{\frac{\lambda}{k+1}}$, $\fixedPoint$ is a fixed point for $\fpMap{}$.

    In addition, we can obtain a closed form for $\cMass{\tIw}$ by dividing
    Equation~\ref{eq:fixed_point_proof_1} by Equation~\ref{eq:fixed_point_proof_2}:
    \begin{equation*}
        \cMass{\tIw} = \frac{\avgPost[\paramValue][\cMass{}]}{\avgPost[\w][\cMass{}]},
    \end{equation*}
    which completes the proof.
\end{proof}

Using Equation~\ref{eq:pus_frontal}, we can construct a sequence of unnormalized measures by
\[\rho_{n+1} = \fpMap{\rho_n} \cdot \]
This sequence can be used in an iterative algorithm to find an estimation of the central distribution
$\cMass{\w}$.
Once $\cMass{\w}$ is determined, Equation~\ref{eq:conditional_for_theta} can be used to obtain $\cMass{\paramValue{}}$.

The key property of Equation~\ref{eq:pus_frontal} is that it allows us to compute $\cMass{\paramValue}$ for
specific regions of the parameter space without needing to scan the entire space, which is typically large and intractable.
Additionally, it defines a mapping based on \w and its marginals \w[i], rather than \paramValue.
Notably, the space of \w is significantly smaller than that of \paramValue.

\begin{definition}[Posterior Potential]
    Let \pmodel be a parametric model and \Qset be a query set.
    Let $\X \in \Qset$, and $\W$ be a function of the parameter \param that fully parametrizes \X.
    We define the \emph{posterior potential} of \X for \w, denoted by $\avgPost[\w][\X]$, as
    follows:
    \begin{equation*}
        \avgPost[\w][\X] = \avgPostExpand[\cMass{}]{\w}{\x} \cdot
    \end{equation*}
    When \X is clear from the context we will omit \X from the notation and simply write $\avgPost[\w]$.
\end{definition}

In Equation~\ref{eq:pus_frontal}, if we replace $\avgPost[][\rho]$ with $\avgPost[][\cMass{}]$, that could improve the
convergence properties considerably.
In a model that the posterior of \X converges, we can find $\avgPost[][\X]$ by using any well behaved prior
instead of $\cMass{\paramValue}$.
We typically will use a conjugate prior to simplify the derivation of the posterior potential.
Using a conjugate prior, the posterior potential takes a form similar to the entropy of the likelihood distribution.

\begin{theorem}
    Let $\fpMap{}$ be the map defined in theorem~\ref{th:fixed_point}.
    The sequence \[\rho_{n+1} = \fpMap{\rho_n},\] converges to the fixed point of
    $\fpMap{}$ for any arbitrary initial unnormalized probability measure $\rho_1$.
\end{theorem}

\begin{proof}
    \renewcommand{\d}[2]{d_{\W} \ifthenelse{\isempty{#1}}{}{\left( #1, #2 \right)}}
    \newcommand{\dExpand}[2]{\intAbs{\lnfrac{#1}{#2}}}
    \newcommand{\lnfrac}[2]{\ln \frac{#1(\w)}{#2(\w)}}
    \newcommand{\intAbs}[1]{\int \left| #1 \right| d\w}
    \newcommand{\p}{\rho}
    \newcommand{\q}{\eta}
    \renewcommand{\r}{\phi}
    We prove the theorem by using Banach fixed point theorem~\cite{Agarwal2018}.
    Let $\p$ and $\q$ be two unnormalized probability measures of a random variable \w
    We define the distance of $\p$ and $\q$ as follows:
    \begin{equation}
        \d{\p}{\q} = \dExpand{\p}{\q} \cdot
        \label{eq:d_metric}
    \end{equation}
    Obviously, if $\p = \q$, then $\d{\p}{\q} = 0$.
    Conversely, if $\d{\p}{\q} = 0$, that implies  $\p = \q$ because $\intAbs{f(\w)} = 0$ if and only if $f(\w) = 0$.
    Moreover, $\d{}{}$ is symmetric $\d{\p}{\q} = \d{\q}{\p}$, and satisfies the triangle inequality:
    \begin{align*}
        \d{\p}{\r} + \d{\r}{\q} &= \dExpand{\p}{\r} + \dExpand{\r}{\q} \\
        &= \int \left( \left| \lnfrac{\p}{\r} \right| + \left| \lnfrac{\r}{\q} \right| \right) d\w \\
        &\geq \intAbs{\lnfrac{\p}{\r} + \lnfrac{\r}{\q}} \\
        &= \d{\p}{\q} \cdot
    \end{align*}
    Thus, $\d{}{}$ is a metric.

    Now, we show that $\fpMap{}$ is a contraction.
    \begin{align*}
        \d{\fpMap{\p}}{\fpMap{\q}} &= \dExpand{\fpMap{\p}}{\fpMap{\q}} \\
        &= \intAbs{ \ln \left( \frac{\p(\w)}{\q(\w)} \pusFrontal[\q][\p]^{\frac{1}{k+1}} \right) }\\
        &= \frac{1}{k+1} \intAbs{ \ln \left( \prod_i \frac{p_{\p}(\w \mid \w[i])}{p_{\q}(\w \mid \w[i])} \right)} \\
        &\leq \frac{1}{k+1} \sum_{i=1}^{k} \intAbs{\ln \frac{p_{\p}(\w \mid \w[i])}{p_{\q}(\w \mid \w[i])}} \\
        &\leq^{?} \frac{k}{k+1} \d{\p}{\q}
    \end{align*}
    Unfortunately, this proof is not complete.
\end{proof}

